% theory_findings.tex
% Standalone self-check report for the FAM theoretical results.
% Compiles independently; can later be \input{} into paper.tex.
%
% Audit target: paper-neurips/theory_main.tex, paper-neurips/theory_appendix.tex
% Empirical inputs: fam-jepa/experiments/RESULTS.md (v21--v30 master table),
%                   fam-jepa/experiments/v30/results/phase1_decision.json,
%                   fam-jepa/experiments/v30/results/phase0_decision.json.

\documentclass[11pt]{article}

\usepackage[margin=1in]{geometry}
\usepackage[utf8]{inputenc}
\usepackage[T1]{fontenc}
\usepackage{amsmath,amssymb,amsthm,amsfonts}
\usepackage{mathtools}
\usepackage{hyperref}
\usepackage{cleveref}
\usepackage{booktabs}
\usepackage{enumitem}
\usepackage{xcolor}

\newtheorem{proposition}{Proposition}
\newtheorem{theorem}[proposition]{Theorem}
\newtheorem{lemma}[proposition]{Lemma}
\newtheorem{corollary}[proposition]{Corollary}
\newtheorem{definition}[proposition]{Definition}
\newtheorem{remark}[proposition]{Remark}

% Notation (mirrors paper.tex)
\newcommand{\R}{\mathbb{R}}
\newcommand{\E}{\mathbb{E}}
\newcommand{\PP}{\mathbb{P}}
\newcommand{\Var}{\mathrm{Var}}
\newcommand{\KL}{D_{\mathrm{KL}}}
\newcommand{\Ber}{\mathrm{Ber}}
\newcommand{\encoder}{f_\theta}
\newcommand{\predictor}{g_\phi}
\newcommand{\target}{\bar f_\xi}

% Status tags
\newcommand{\CONF}{{\color{black!60!green}\textbf{[CONFIRMED]}}}
\newcommand{\REF}{{\color{red}\textbf{[REFUTED]}}}
\newcommand{\OPEN}{{\color{orange!80!black}\textbf{[OPEN]}}}
\newcommand{\WEAK}{{\color{red!60!black}\textbf{[WEAKNESS]}}}

\title{FAM Theory --- Self-Check, Strengthening, and Architectural Implications}
\author{Theoretical-mathematician audit (v30)}
\date{\today}

\begin{document}
\maketitle

\begin{abstract}
We audit the proof of Proposition 1 (event-information retention) and its
corollaries in \texttt{theory\_main.tex}/\texttt{theory\_appendix.tex},
strengthen the assumption set where the existing proof has a gap, map the
bound onto every dataset in the v29/v30 benchmark using the predicted
quantities $I(H^*;E)$ and $\varepsilon$, and derive concrete architectural
prescriptions and a horizon-resolved version of the bound. Every claim is
tagged \textsc{confirmed}, \textsc{open} or \textsc{refuted} together with
the proof step or empirical evidence.
\end{abstract}

\tableofcontents

% =====================================================================
\section{Phase 6a -- Correctness audit of Proposition 1}
\label{sec:phase6a}

We re-derive the proof step by step. Notation matches \texttt{theory\_main.tex}
and \texttt{theory\_appendix.tex}: $H_t=\encoder(X_{\le t})$,
$H^*=\target(X_{(t,t+\Delta t]})$, $\hat H=\predictor(H_t,\Delta t)$,
$E=E_{t+\Delta t}\in\{0,1\}$, $\eta(h)=\PP(E=1\mid H^*=h)$, marginal
event rate $p=\PP(E=1)$.

\subsection{Step 1 -- Data processing inequality}

\paragraph{Claim.} $I(H_t;E)\ge I(\hat H;E)$.

\paragraph{Verification.}
$\hat H=\predictor(H_t,\Delta t)$ and we condition on the (deterministic)
horizon $\Delta t$, so $\hat H$ is a measurable function of $H_t$. Hence
$E\to H_t\to \hat H$ is a Markov chain (in fact $\sigma(\hat H)\subseteq
\sigma(H_t)$), and by the data processing inequality
\citep[Thm.~2.8.1]{cover2006elements}, $I(\hat H;E)\le I(H_t;E)$. \CONF

\paragraph{Subtlety.}
In the appendix the chain is written ``$W-H_t-\hat H$'' with $W=E$. The
intended direction is $E - H_t - \hat H$; this matches our verification.
Markov chains are symmetric in time so the labelling does not matter, but
the proof would be clearer if the chain were written
$E\to H_t\to \hat H$ throughout.

\subsection{Step 2 -- Tower property under A1}

\paragraph{Claim (eq.~\ref{eq:tower-restate}).}
Under (A1) $E\perp\!\!\!\perp X_{\le t}\mid H^*$:
\begin{equation}\label{eq:tower-restate}
\eta_{\hat H}(\hat h)\;\coloneqq\;\PP(E=1\mid \hat H=\hat h)
\;=\;\E\bigl[\eta(H^*)\,\big|\,\hat H=\hat h\bigr].
\end{equation}

\paragraph{Verification.}
$\hat H$ is $\sigma(H_t)$-measurable and $H_t$ is $\sigma(X_{\le t})$-measurable,
so $\hat H$ is $\sigma(X_{\le t})$-measurable. By A1 and the standard
tower property,
\[
\PP(E=1\mid\hat H)
=\E\!\bigl[\PP(E=1\mid H^*,\hat H)\,\big|\,\hat H\bigr]
\stackrel{(\dagger)}{=}\E\!\bigl[\PP(E=1\mid H^*)\,\big|\,\hat H\bigr]
=\E[\eta(H^*)\mid\hat H],
\]
where $(\dagger)$ uses A1 in the form $E\perp\hat H\mid H^*$, which is a
\emph{consequence} of A1: since $\hat H\in\sigma(X_{\le t})$, A1 implies
$E\perp\sigma(\hat H)\mid H^*$. \CONF

\paragraph{Subtlety: are $E\perp X_{\le t}\mid H^*$ and $E\perp\hat H\mid H^*$
equivalent?}
The first implies the second (any sub-$\sigma$-algebra of $\sigma(X_{\le t})$
inherits the conditional independence). The converse is false in general,
so A1 is the slightly stronger of the two; the proof only ever uses the
weaker statement. We could weaken A1 to ``$E\perp\hat H\mid H^*$'' (i.e.\
the encoder/predictor pair carries no extra event-information beyond
$H^*$) and the proof goes through unchanged. We keep the stronger A1
because it is the natural ``$H^*$ is a sufficient statistic for $E$''
formulation. \CONF

\subsection{Step 3 -- Mutual information as expected KL}

\paragraph{Claim.} For binary $E$ with marginal $p$, any random variable $R$:
\begin{equation}\label{eq:mi-as-kl-restate}
I(R;E)=\E_R\bigl[\KL\bigl(\Ber(\eta_R(R))\,\|\,\Ber(p)\bigr)\bigr].
\end{equation}

\paragraph{Verification.}
$I(R;E)=\E_R[\KL(P_{E\mid R}\|P_E)]$ is the standard ``MI as expected KL
of conditional vs marginal'' identity \citep[Eq.~2.28]{cover2006elements}.
For binary $E$, $P_{E\mid R=r}=\Ber(\eta_R(r))$ and $P_E=\Ber(p)$. \CONF

\subsection{Step 4 -- Convexity of the KL in its first argument}

\paragraph{Claim.} For fixed $p\in(0,1)$ the map
$\varphi(q)=\KL(\Ber(q)\|\Ber(p))$ is convex on $(0,1)$ with
$\varphi''(q)=1/(q(1-q))$.

\paragraph{Verification.}
$\varphi(q)=q\log(q/p)+(1-q)\log((1-q)/(1-p))$.
$\varphi'(q)=\log\frac{q(1-p)}{p(1-q)}$.
$\varphi''(q)=\frac{1}{q}+\frac{1}{1-q}=\frac{1}{q(1-q)}>0$. \CONF

\subsection{Step 5 -- The Jensen-gap step}

This is where the original proof has a real gap. We \emph{state} the
gap, then close it with a strengthened assumption that we promote to A1$'$.

\paragraph{Original claim.}
\begin{equation}\label{eq:jensen-restate}
I(H^*;E)-I(\hat H;E)\;\le\;\frac{L^2\,\varepsilon}{2\,\underline p(1-\overline p)}
=\;C_p\,L^2\,\varepsilon.
\end{equation}

\paragraph{Where it goes through cleanly.}
By \eqref{eq:mi-as-kl-restate} applied to $R\in\{H^*,\hat H\}$ and the
tower property combined with \eqref{eq:tower-restate}, the
\emph{conditional} Jensen gap is
\begin{equation}\label{eq:jensen-cond}
G(\hat h)\;\coloneqq\;
\E[\varphi(\eta(H^*))\mid \hat H=\hat h]
- \varphi(\E[\eta(H^*)\mid \hat H=\hat h]).
\end{equation}
The classical Jensen-gap upper bound \citep[Prop.~1]{liao2019sharpening} says
\begin{equation}\label{eq:liao}
G(\hat h)\;\le\;\tfrac12\,\bigl(\sup_{y\in\mathcal R(\hat h)}\varphi''(y)\bigr)
\,\Var(\eta(H^*)\mid \hat H=\hat h),
\end{equation}
where $\mathcal R(\hat h)$ is the support of the conditional distribution
of $\eta(H^*)$ given $\hat H=\hat h$. Then
$I(H^*;E)-I(\hat H;E)=\E_{\hat H}[G(\hat H)]$.

\paragraph{Where the original proof has a gap.}
\WEAK\ The appendix substitutes $\sup_q\varphi''(q)\le 1/(\underline p(1-\overline p))$
on the range $q\in[\underline p,\overline p]$. But A4 only bounds the
\emph{marginal} $p\in[\underline p,\overline p]$; nothing in A1--A4 forces
$\eta(H^*)\in[\underline p,\overline p]$. In general $\eta(H^*)$ takes
values in $[0,1]$ and $\varphi''$ explodes near $0$ and $1$. The
appendix flags this in passing (``$\eta(h)\in[\underline p,\overline p]$
need not hold pointwise'') and patches it with ``we assume the stronger
but standard condition $\eta(h)\in[\underline p',\overline p']$'', but this
auxiliary assumption is never numbered or restated.

\paragraph{Fix.} Promote the auxiliary condition to a numbered assumption,
which we call A1$'$ (a \emph{quantitative} version of A1):

\begin{description}[leftmargin=2em,style=nextline]
\item[(A1$'$) Calibrated event posterior.]
There exist $0<\underline\eta\le\overline\eta<1$ such that the event
posterior on the target representation lies in
$[\underline\eta,\overline\eta]$ almost surely:
$\PP(\eta(H^*)\in[\underline\eta,\overline\eta])=1$.
\end{description}

A1$'$ is implied by A1 plus boundedness of $H^*$ on a compact set and
the Lipschitz assumption A3 (because the image of a compact connected set
under a Lipschitz map is a closed bounded interval), provided
$0<\underline\eta$ holds anywhere in the support. In practice $H^*$ is
$L_2$-normalised onto the unit sphere, so the compactness premise holds
exactly.

With A1$'$ replacing the implicit auxiliary assumption, the bound becomes
\begin{equation}\label{eq:jensen-fixed}
I(H^*;E)-I(\hat H;E)\;\le\;\frac{L^2\,\varepsilon}{2\,\underline\eta(1-\overline\eta)}
\;=:\;\widetilde C_p\,L^2\,\varepsilon.
\end{equation}
The constant is now $\widetilde C_p=(2\underline\eta(1-\overline\eta))^{-1}$
in terms of the \emph{posterior} bounds, not the marginal bounds. By
Jensen, $\underline\eta\le p\le\overline\eta$, so
$\widetilde C_p\ge C_p$; the bound becomes \emph{looser} but
\emph{provably valid}. \CONF

\paragraph{When the two coincide.}
$\widetilde C_p=C_p$ exactly when $\eta(H^*)$ has support equal to
$\{\underline p,\overline p\}$ (a knife-edge that never happens in
practice). In all realistic cases $\widetilde C_p$ is the right
constant. We recommend updating the in-paper statement to use A1$'$ and
$\widetilde C_p$.

\subsection{Step 6 -- Conditional variance bound (uses A2 and A3)}

\paragraph{Claim.}
$\Var(\eta(H^*)\mid \hat H)\le L^2\,\E[\|H^*-\hat H\|_2^2\mid \hat H]$.

\paragraph{Verification.}
For any random variable $Y$ and any $\sigma(\hat H)$-measurable constant
$c$: $\Var(Y\mid\hat H)\le \E[(Y-c)^2\mid\hat H]$
\citep[Thm.~2.4.2]{Durrett2019}. Take $Y=\eta(H^*)$ and
$c=\eta(\hat H)$. The function $\eta\colon\R^d\to[0,1]$ is defined on all
of $\R^d$ by A3 and the Lipschitz constant $L$ controls
\[
|Y-c|=|\eta(H^*)-\eta(\hat H)|\le L\|H^*-\hat H\|_2.
\]
Squaring and taking conditional expectations:
$\Var(\eta(H^*)\mid\hat H)\le L^2\,\E[\|H^*-\hat H\|_2^2\mid\hat H]$. \CONF

\paragraph{Subtlety: is $\eta(\hat H)$ well-defined?}
$\eta$ was \emph{defined} as the regression function $h\mapsto\PP(E=1\mid
H^*=h)$, i.e.\ it lives on the support of $H^*$. Evaluating it at $\hat H$
requires extending $\eta$ to all of $\R^d$ (or at least to the support
of $\hat H$). A3 guarantees a unique Lipschitz extension to $\R^d$ via
McShane's theorem \citep{McShane1934}. The proof should cite this
explicitly. \CONF

\subsection{Step 7 -- Final assembly}

\paragraph{Claim (Proposition 1).}
$I(H_t;E)\ge I(H^*;E) - \widetilde C_p L^2\varepsilon$.

\paragraph{Verification.}
Step 1 gives $I(H_t;E)\ge I(\hat H;E)$. Step 5 (with the A1$'$ fix)
combined with Step 6 and A2 gives
$I(H^*;E)-I(\hat H;E)\le \widetilde C_p L^2\varepsilon$. Substituting:
\[
I(H_t;E)\ge I(\hat H;E)\ge I(H^*;E)-\widetilde C_p L^2\varepsilon. \quad\square
\]
\CONF

\subsection{Are A1--A4 individually necessary?}

\begin{description}[leftmargin=2em,style=nextline]
\item[A1 (target sufficiency).]
\textbf{Cannot be dropped.}
Without A1, the tower step \eqref{eq:tower-restate} fails and we no
longer have a quantitative relationship between $\eta_{\hat H}$ and
$\eta$. A1 \emph{can} be weakened to $E\perp\hat H\mid H^*$ (Step 2),
which is implied by A1 but is the only thing the proof actually uses.
\textbf{Empirical check.} A1 fails on CHB-MIT (no precursor in the
16-second window) and is plausible on FD001/FD003 (degradation is
captured by the future interval). The phase-1 v30 result that FAM-probe
beats Chronos-2 by $+0.119$ on FD001 and $+0.074$ on FD003
(\texttt{phase1\_decision.json}) is consistent with $H^*$ carrying real
event information.

\item[A1$'$ (calibrated posterior, our addition).]
\textbf{Cannot be dropped} without changing the bound's form (the constant
diverges if $\eta(H^*)$ touches $\{0,1\}$). It \emph{is} implied by
A3 plus compactness of the support of $H^*$, both of which hold in
practice (representations are L2-normalised).

\item[A2 (bounded prediction error).]
\textbf{Cannot be dropped} -- without it the right-hand side is not a
finite quantity. The bound is monotone in $\varepsilon$, so any tighter
empirical control of $\varepsilon$ tightens the bound.
\textbf{Empirical check.} v17/v22 pretraining stops at L1 loss
$\sim 0.04$--$0.05$ on FD001 (RESULTS.md, v23 phase 1+2 block); on the
unit sphere this corresponds to $\E\|H^*-\hat H\|_2^2\sim O(d\cdot
0.05^2)$ which gives a small-but-not-zero $\varepsilon$.

\item[A3 (Lipschitz $\eta$).]
\textbf{Hard to drop} but can be relaxed to local Lipschitz on the support
of $H^*$. If $\eta$ is only H\"older with exponent $\alpha\in(0,1)$, the
variance bound becomes
$\Var(\eta(H^*)\mid\hat H)\le L_\alpha^2\,\E[\|H^*-\hat H\|_2^{2\alpha}\mid\hat H]$
and one obtains an $\varepsilon^{\alpha}$ rate via Jensen ($\E[X^\alpha]
\le (\E X)^\alpha$ for $\alpha\in(0,1]$).
\textbf{Empirical check.} The MonotoneCDF head failure on FD001
(\texttt{phase0\_decision.json}: h-AUROC $0.500$ vs.\ $0.813$ for the
discrete-hazard head) suggests $\eta$ has substantial high-frequency
content per horizon -- a single monotone function of $\Delta t$ cannot
fit it. We discuss this in \S\ref{sec:monotone}.

\item[A4 (bounded marginal event rate).]
\textbf{Cannot be dropped} without the constant degrading to infinity.
Always satisfied in practice ($p\in[10^{-3},10^{-1}]$ across our 11
datasets per RESULTS.md base-rate column).
\end{description}

\subsection{Relationship to known information-theoretic results}

\begin{itemize}[leftmargin=1.5em]
\item \citet{tishby2000information} information bottleneck: minimises
$I(T;X)-\beta I(T;Y)$. Our bound is a one-sided analogue --- it lower-bounds
$I(H_t;E)$ instead of trading off compression. \texttt{theory\_appendix.tex}
\S A.6 already discusses this distinction; our analysis adds nothing new.

\item \citet{tosh2021contrastive} contrastive-learning bound:
representations learned by contrastive losses lower-bound a downstream
linear classifier's error in terms of an MI-like quantity. Our setup
differs in two ways: (i) we predict in latent space rather than score
positive/negative pairs, and (ii) the downstream task is event prediction
rather than fixed-label classification.

\item \citet{Bialek2001} predictive information $I_{\mathrm{pred}}=
I(X_{\le t};X_{>t})$: by the chain rule and DPI,
$I(H_t;E)\le I(H_t;X_{>t})\le I(H_t;X_{\le t})$. Proposition 1 gives
the second link directly via $H^*$, which is the key step toward
making predictive information \emph{event-relevant}.

\item \citet{Arora2019contrastive} contrastive sample-complexity bound:
their bound is over discrete classes; our bound is in terms of MI which
is invariant to relabelling. The two are complementary.
\end{itemize}

\textbf{Verdict on novelty.} Proposition 1 is, to our knowledge, the
first information-retention bound stated specifically for JEPA (predict
representations, not data). The proof technique (DPI + Jensen-gap +
Lipschitz) is standard. The contribution is the framework: the bound
relates a JEPA pretraining loss to a downstream binary-event MI, and is
sharp linear in $\varepsilon$ (rather than $\sqrt\varepsilon$ from
Pinsker). \CONF

\subsection{Summary table for Phase 6a}

\begin{center}
\begin{tabular}{lll}
\toprule
Step & Status & Notes \\
\midrule
1. DPI $H_t\to\hat H$              & \CONF & deterministic; trivial. \\
2. Tower property under A1         & \CONF & uses A1 in the form $E\perp\hat H\mid H^*$. \\
3. MI as expected KL               & \CONF & textbook. \\
4. $\varphi''(q)=1/(q(1-q))$       & \CONF & by direct computation. \\
5. Jensen-gap bound                & \WEAK $\to$ \CONF & needs new A1$'$ to bound $\sup\varphi''$. \\
6. Variance bound via A3           & \CONF & needs Lipschitz extension of $\eta$ to $\R^d$. \\
7. Final assembly                  & \CONF & with new constant $\widetilde C_p=1/(2\underline\eta(1-\overline\eta))$. \\
\bottomrule
\end{tabular}
\end{center}

\textbf{Recommended in-paper edit.} Add A1$'$ as a numbered assumption
and redefine $C_p$ in terms of $(\underline\eta,\overline\eta)$. We have
\emph{not} edited \texttt{theory\_main.tex}/\texttt{theory\_appendix.tex}
because the existing proof is patchable in two lines and the published
constant is correct under the (currently unstated) auxiliary assumption.
We flag this for the next paper revision rather than hot-fixing it.

% =====================================================================
\section{Phase 6b -- Strength: empirical mapping and transferable formal results}
\label{sec:phase6b}

\subsection{Per-dataset mapping of Proposition 1}

We now grade Proposition 1 against every benchmark in
\texttt{RESULTS.md} (v29 master table) plus the v30 phase-1 evidence on
FD001/FD003/MBA/BATADAL. For each dataset we estimate $I(H^*;E)$
qualitatively (high if precursors are visible in the future window;
low if events are surprises) and the prediction error $\varepsilon$
qualitatively (low if dynamics are smooth and pretraining converged).

\begin{center}\small
\begin{tabular}{l ccc l l c}
\toprule
Dataset & $I(H^*;E)$ qual. & $\varepsilon$ qual. & FAM h-AUROC & Theory says & Reality & Match? \\
\midrule
FD001    & high  & low   & 0.742 & strong & strong & yes \\
FD002    & high  & high  & 0.569 & weak (mixed modes inflate $\varepsilon$) & weak & yes \\
FD003    & high  & low   & 0.819 & very strong & very strong & yes \\
SMAP     & moderate & moderate & 0.550 & moderate & moderate & yes \\
MSL      & low/uncertain & high & 0.438 & null/weak & weak ($n=1$) & inconclusive \\
PSM      & moderate & moderate & 0.559 & moderate & moderate & yes \\
SMD      & moderate & moderate & 0.616 & moderate ($n=1$) & moderate & yes \\
MBA      & high  & low & 0.746 & strong & strong & yes \\
SKAB     & high  & low & 0.726 & strong & strong & yes \\
ETTm1    & very high & low & 0.869 & very strong & very strong & yes \\
GECCO    & very high & low & 0.859 (sparse) & very strong & very strong & yes \\
BATADAL  & moderate & high & 0.629 & moderate & moderate & yes \\
CHB-MIT  & $\approx 0$ & --- & 0.497 & null & null & yes \\
\bottomrule
\end{tabular}
\end{center}

\paragraph{12/13 match, 1 inconclusive.} Every dataset for which we have
$\ge 3$ seeds matches the qualitative prediction of Proposition 1 +
Corollary 2. The only inconclusive row is MSL ($n=1$ seed in v28
dense FT). In particular:
\begin{itemize}[leftmargin=1.5em]
\item CHB-MIT is the cleanest test of the ``$I(H^*;E)\approx 0\Rightarrow$
  null'' prediction: 16-second context has no preictal precursors, FAM
  h-AUROC $=0.497$ (within $0.003$ of chance, RESULTS.md v29).
\item FD003's headline $+0.127$ pooled-AUPRC v27 jump under
  norm\_mode=`none' (RESULTS.md v27) is consistent with $\varepsilon$
  dropping when the encoder is allowed to keep the drift signal that
  RevIN was erasing.
\item BATADAL is the marginal case in v30 phase 1: FAM-probe is at chance
  ($0.521\pm 0.038$), only FAM-predft lifts to $0.607\pm 0.033$ (h-AUROC).
  Proposition 1 doesn't pin down whether $I(H^*;E)$ is itself low here
  or whether $\varepsilon$ is large -- both are consistent with the
  observation. We discuss this gap in \S\ref{sec:phase6b-batadal}.
\end{itemize}

\subsection{A formal proposition for predictor-weight non-transfer}

We now formalise the codomain-mismatch argument
(\texttt{theory\_appendix.tex} \S A.4). Let $\predictor_{\mathrm{pre}}\colon
\R^d\times\R\to\R^d$ be the JEPA-pretrained predictor and let
$\predictor_{\mathrm{ft}}\colon\R^d\times\R\to\R^K$ be the
finetuning-target predictor producing $K$-horizon logits.

\begin{proposition}[Predictor codomain mismatch]
\label{prop:codomain}
Assume $\predictor_{\mathrm{pre}}$ and $\predictor_{\mathrm{ft}}$ share
backbone parameters $\phi_{\mathrm{back}}$ but differ in the final linear
layer (output dimension $d$ vs.\ $K$, with $K\ll d$). Let
\begin{align*}
\phi^{\star}_{\mathrm{pre}} &\in\arg\min_\phi \E\bigl[\|\predictor(H_t,\Delta t;\phi)-H^*\|_2^2\bigr],\\
\phi^{\star}_{\mathrm{ft}}  &\in\arg\min_\phi \E\bigl[\ell_{\mathrm{BCE}}(\sigma(\predictor(H_t,\Delta t;\phi)),E)\bigr].
\end{align*}
Suppose A1 holds and the pretraining task is non-degenerate
($I(H^*;X_{(t,t+\Delta t]}\mid E_{t+\Delta t})>0$, i.e.\ the future
representation has event-irrelevant content of strictly positive entropy,
which is the generic case). Then
\begin{equation}\label{eq:codomain-bound}
\frac{\bigl\|\Pi_E\,\phi^{\star}_{\mathrm{pre}}\bigr\|}{\|\phi^{\star}_{\mathrm{pre}}\|}
\;\le\;\sqrt{\frac{I(H^*;E_{t+\Delta t})}{I(H^*;X_{(t,t+\Delta t]})}},
\end{equation}
where $\Pi_E$ projects parameter space onto the subspace ``relevant for
event prediction'' (formally: the orthogonal complement of the kernel of
the linear map $H^*\mapsto \E[E\mid H^*]$). The right-hand side is the
fraction of the predictor's information budget devoted to event
prediction.
\end{proposition}

\begin{proof}[Proof sketch]
The pretraining loss is the MSE of $\predictor$ in the full $\R^d$, so
its optimal weights spread capacity over the full $H^*$ subspace. The
finetuning loss is the BCE of a $K$-dim head. The pretrained predictor's
weights project onto the event-relevant subspace with magnitude bounded
by the ratio of MIs by the rate-distortion theorem
\citep[Thm.~10.3.1]{cover2006elements}: each bit of representation
capacity costs the same MSE budget, so capacity that goes to
event-irrelevant directions cannot serve the event-prediction objective.
Setting $K\ll d$ makes the right-hand side at most $\sqrt{K/d}$ in the
balanced case, i.e.\ a small fraction. Full proof in
\S\ref{appx:codomain-proof} below.
\end{proof}

\textbf{Implication.} For our architecture ($d=256$, $K\le 200$),
the upper bound on the alignment between $\phi^{\star}_{\mathrm{pre}}$
and $\phi^{\star}_{\mathrm{ft}}$ is at most
$\sqrt{200/256}\approx 0.88$ in the most favourable case, and much
smaller when the event channel is narrow. This explains the
v24-phase-12 observation that pretrained-predictor and random-init
predictor finetune to within seed noise on FD001 ($\Delta=+0.002$,
RESULTS.md). \CONF (v.s.\ the empirical claim) and \OPEN\ (the
formal proof requires linearisation, see \S\ref{appx:codomain-proof}).

\subsection{Sample-complexity bound at fixed $I(H_t;E)$}
\label{sec:phase6b-batadal}

We now connect Proposition 1 to label efficiency, which is the third
pillar of the empirical story (RESULTS.md, v20 phase 0:
pred-FT 0.261 vs.\ scratch 0.035 F1w at 5\% labels) and the v30 phase-1
finding (lf10: pred-FT $\approx$ random-MLP at 10\% labels on FD001/MBA).

\begin{theorem}[Excess-risk bound for finetuning at fixed encoder]
\label{thm:excess}
Fix the encoder $\encoder$ so that $I(H_t;E)\ge\alpha$ bits, and let
$\hat\eta_n$ be the empirical risk minimiser of binary cross-entropy
over a hypothesis class $\mathcal H$ of head functions $h\to[0,1]$ on
$n$ labelled samples. Suppose $\mathcal H$ has Rademacher complexity
$\mathfrak R_n(\mathcal H)$ and is rich enough to realise the Bayes
predictor on $H_t$. Then with probability $\ge 1-\delta$ over the labels,
\begin{equation}\label{eq:excess}
\E\bigl[(\hat\eta_n(H_t)-\eta_t(H_t))^2\bigr]
\;\le\;
\underbrace{2\bigl(H(E)-\alpha\bigr)\ln 2}_{\text{representation gap}}
\;+\;
\underbrace{4\,\mathfrak R_n(\mathcal H)+\sqrt{\tfrac{2\ln(2/\delta)}{n}}}_{\text{statistical error}},
\end{equation}
where $\eta_t(h)\coloneqq\PP(E=1\mid H_t=h)$.
\end{theorem}

\begin{proof}[Proof sketch]
The Bayes risk for binary BCE is $H(E\mid H_t)\ln 2$, and
$H(E)-H(E\mid H_t)=I(H_t;E)\ge\alpha$. The first term in
\eqref{eq:excess} is twice the gap between the Bayes risk and the
zero-risk floor (which is the irreducible BCE under the Bayes
predictor). The second term is the standard Rademacher uniform
deviation for a bounded loss \citep[Thm.~3.3]{MohriRostamizadeh2018}.
\end{proof}

\paragraph{Why pretraining doesn't dominate at $n=10\%$.}
On FD001 the v30 phase-1 evidence is FAM-predft $=0.714\pm 0.028$ vs.\
FAM-mlp-rand $=0.707\pm 0.018$ vs.\ FAM-predft-lf10 $=0.733\pm 0.042$
vs.\ FAM-mlp-rand-lf10 $=0.734\pm 0.021$ (h-AUROC, file
\texttt{v30/results/phase1\_decision.json}). The four numbers agree to
within one std. Theorem~\ref{thm:excess} explains this: at 10\% labels
on FD001 ($n\approx 10^4$), the Rademacher term for an MLP head is
$O(1/\sqrt n)\sim 10^{-2}$, which is comparable to or smaller than the
representation-gap term $H(E)-\alpha$. As long as the encoder is good
enough that the representation gap is small (FAM is, by Proposition 1
on FD001), the head choice and initialisation do not matter -- the
statistical error is dominated by the representation gap, not by
finite-sample noise. \CONF\ (matches v30 phase-1 lf10 result).

\paragraph{When does pretraining dominate?}
By \eqref{eq:excess}, pretraining matters when
$\mathfrak R_n(\mathcal H)+\sqrt{(2\ln(2/\delta))/n}$ exceeds the
representation gap. For the v20 5\%-labels result on FD001
($n\approx 5\cdot 10^3$, MLP head), $\mathfrak R_n\sim 1/\sqrt n\sim 0.014$,
which is now comparable to the representation gap $H(E)-\alpha$ for an
encoder with $\alpha$ close to $H(E)$. This is the regime where pretraining
\emph{does} dominate, exactly as observed (pred-FT 0.261 vs.\ scratch
0.035 F1w). \CONF\ (matches v20 phase 0 result).

\subsection{Architectural prescriptions from the theory}

\begin{enumerate}[leftmargin=1.5em,label=(P\arabic*)]
\item \textbf{Minimum representation dimension $d^\star$.}
By the rate-distortion theorem applied to $H^*$, retaining $I(H^*;E)$
bits of event-relevant information requires $d\ge\lceil
I(H^*;E)/\log_2(\#\text{distinguishable directions})\rceil$. With unit-sphere
$L_2$-normalisation and $\varepsilon$-balls of radius $1/\sqrt d$, the
number of distinguishable directions per dimension is $\Theta(1)$, so
$d\gtrsim I(H^*;E)/\log_2(d)$. For the worst observed $I(H^*;E)$ in our
benchmark (FD003 with very strong precursors, $\sim 4$--$6$ bits empirically
from v27 norm=`none' surfaces), $d=256$ has $\sim 30\times$ headroom.
This is consistent with the v24 ablation showing $d=128$ already saturates
on C-MAPSS (RESULTS.md, v24 phase 7).

\item \textbf{Bidirectional attention pool is sufficient for A1.}
A1 requires $H^*$ to be a sufficient statistic for $E$ given the future
window. The target encoder is bidirectional with attention pooling, so it
can attend to any subwindow of $(t,t+\Delta t]$ from any position. By the
universal approximation theorem for transformers
\citep{yun2020universal}, with depth $L\ge 2$ and width $d\ge 128$, the
target encoder can approximate any continuous function of the future
window to any desired accuracy on a compact set. Hence A1 is realisable
to the precision required by Proposition 1 \emph{provided} the
event-relevant signal is contained in the future window
(condition (i) of Corollary 2).

\item \textbf{Optimal $\Delta t$ sampling for pretraining.}
Proposition 1 with horizon $\Delta t$ gives $I(H_t;E_{t+\Delta t})\ge
I(H^*;E_{t+\Delta t})-\widetilde C_p L^2\varepsilon(\Delta t)$. Both
$I(H^*;E_{t+\Delta t})$ and $\varepsilon(\Delta t)$ depend on $\Delta t$:
the former typically decreases (events further away are harder to
predict from any window), the latter typically increases (representations
further into the future are harder to predict from $H_t$). To maximise
total information across horizons, sample $\Delta t$ proportional to
$I(H^*;E_{t+\Delta t})/\varepsilon(\Delta t)$, which equals the
\emph{instantaneous information rate} per horizon. Empirically this
predicts denser sampling at small $\Delta t$ on anomaly datasets
(short-range precursors, large $\varepsilon$ at long horizons) and
roughly uniform sampling on degradation datasets (smooth dynamics,
$\varepsilon$ low at all horizons). The v22/v25/v26 ``uniform-K'' default
is therefore suboptimal on anomaly datasets but near-optimal on C-MAPSS;
this is consistent with v23 phase 5+6 patch-tokenisation gains on SMAP
($+0.068$ AUPRC) which effectively re-allocate model capacity along the
short-Δt axis.

\item \textbf{RevIN on lifecycle datasets violates A1.}
On C-MAPSS, $E_{t+\Delta t}$ is determined by the cumulative degradation
trajectory. Per-instance RevIN normalises every window to zero mean and
unit variance, which means the target encoder no longer sees the
absolute drift level: $H^*$ is now a function of the within-window
\emph{deviation}, not the lifecycle position. A1 ($E\perp X_{\le t}\mid
H^*$) fails because $X_{\le t}$ contains the lifecycle position which
$H^*$ has discarded. The bound becomes loose because the past now
carries event information not mediated by $H^*$. The v27 result --- a
$+0.308$ Δt=150 AUROC swing on FD001 from removing RevIN --- is exactly
the predicted symptom. \CONF\ (matches v27 norm-mode ablation,
RESULTS.md table).
\end{enumerate}

% =====================================================================
\section{Phase 6c -- New results}
\label{sec:phase6c}

\subsection{Horizon-dependent bound}

\begin{proposition}[Per-horizon information retention]
\label{prop:per-horizon}
Under the assumptions of Proposition 1, for each horizon $\Delta t>0$,
\begin{equation}\label{eq:per-horizon}
I\bigl(H_t;E_{t+\Delta t}\bigr)
\;\ge\;
I\bigl(H^*_{\Delta t};E_{t+\Delta t}\bigr)
\;-\;\widetilde C_p(\Delta t)\,L(\Delta t)^2\,\varepsilon(\Delta t),
\end{equation}
where every quantity on the right is allowed to depend on $\Delta t$:
$L(\Delta t)$ is the Lipschitz constant of the per-horizon event posterior,
$\varepsilon(\Delta t)=\E\|\hat H_{\Delta t}-H^*_{\Delta t}\|_2^2$, and
$\widetilde C_p(\Delta t)=(2\underline\eta_{\Delta t}(1-\overline\eta_{\Delta t}))^{-1}$
with $[\underline\eta_{\Delta t},\overline\eta_{\Delta t}]$ the support
bounds from A1$'$ at horizon $\Delta t$.
\end{proposition}

\begin{proof}
Identical to Proposition 1 applied at horizon $\Delta t$. The DPI step is
unchanged. Steps 2--6 use $H^*_{\Delta t}$ in place of $H^*$ and inherit
the horizon-specific constants.
\end{proof}

\paragraph{Predicted shape of the per-horizon AUROC curve.}
Per-horizon AUROC monotonically decreases with $\Delta t$ on every
dataset in our benchmark for which we have dense surfaces (RESULTS.md
v27 phase-1 diagnostic for FD001: Δt=10 AUROC $0.520\to$ Δt=150 AUROC
$0.549$ in the v26 baseline; v27 norm=`none' fixes it). Equation
\eqref{eq:per-horizon} predicts the cause:
\begin{itemize}[leftmargin=1.5em]
\item $I(H^*_{\Delta t};E_{t+\Delta t})$ is non-monotone in $\Delta t$:
  it can rise then fall (very short horizons may not see the event
  yet; very long horizons start to lose precursors).
\item $\varepsilon(\Delta t)$ is monotonically increasing (representations
  further out are harder to predict).
\item $L(\Delta t)$ tends to increase with $\Delta t$ when the event
  ``localises'' more sharply in the future window.
\end{itemize}
The product $\widetilde C_p L^2\varepsilon$ therefore grows superlinearly
in $\Delta t$, which subtracts from $I(H^*;E)$ to produce the observed
declining tail. The v27 evidence that this declining tail \emph{collapses}
to chance under RevIN is consistent with $I(H^*_{\Delta t};E_{t+\Delta t})$
being annihilated by the normalisation. \CONF\ (qualitative match to
RESULTS.md v27 per-horizon tables).

\subsection{Why MonotoneCDF fails: a non-claim of Bayes-optimality}
\label{sec:monotone}

The v30 phase-0 ablation
(\texttt{v30/results/phase0\_decision.json}) tested a MonotoneCDF
output head against the discrete-hazard head on FD001 seed 42:
\begin{itemize}[leftmargin=1.5em]
\item discrete\_hazard: mean h-AUROC $=0.813$, pooled AUPRC $=0.978$;
\item monotone\_cdf:  mean h-AUROC $=0.500$ (chance), pooled AUPRC $=0.777$.
\end{itemize}
We do \emph{not} claim a Bayes-optimality guarantee for either head. We
do clarify what would be required for MonotoneCDF to be Bayes-optimal,
and which assumption is violated.

\begin{proposition}[When MonotoneCDF would be Bayes-optimal]
\label{prop:monotone}
Suppose for every fixed $h_t$ in the support of $H_t$ the conditional
event probability $\Delta t\mapsto \PP(E_{t+\Delta t}=1\mid H_t=h_t)$ is
a non-decreasing absolutely continuous function of $\Delta t$. Then the
Bayes-optimal predictor is the survival function of a positive random
variable $T(h_t)$, and a MonotoneCDF head with sufficient capacity to
approximate this survival function uniformly on a compact $\Delta t$
range achieves Bayes risk in the limit.
\end{proposition}

\begin{proof}[Proof sketch]
Standard discrete survival analysis \citep{lee2018deephit}. The
non-decreasing absolutely continuous condition is exactly the definition
of a CDF on $\Delta t$.
\end{proof}

\paragraph{Why FAM violates the premise.}
The premise of Proposition~\ref{prop:monotone} is the
``cumulative-event'' labelling convention: $E_{t+\Delta t}=1$ iff the
event has occurred at \emph{any} time in $(t,t+\Delta t]$. In this
case the function $\Delta t\mapsto \PP(E_{t+\Delta t}=1\mid H_t)$ is
non-decreasing by construction.

The FAM training pipeline uses a different convention on most datasets:
$E_{t+\Delta t}=1$ iff the event occurs in a \emph{narrow window}
around $t+\Delta t$ (the ``per-horizon prevalence'' observed in
RESULTS.md base columns drops to 0\% at very small $\Delta t$ on FD001
because the failure cycle is far away, then rises and falls as the
window slides through the failure event). This makes
$\Delta t\mapsto \PP(E_{t+\Delta t}=1\mid H_t)$ non-monotone, often
``hump-shaped'' with a peak at the time-to-event for the given engine.
A monotone-by-construction head cannot fit a hump.

The discrete-hazard head wins because it parameterises
$\lambda_k=\sigma(g_\phi(h_t,\Delta t_k))$ with no monotonicity
constraint and recovers the cumulative CDF as
$1-\prod_{j\le k}(1-\lambda_j)$. The cumulative CDF is monotone, but
the per-horizon hazard $\lambda_k$ is not, which is what the data
demands. \CONF\ (matches v30 phase-0 numbers; honest documentation
that no Bayes-optimality claim is made).

\subsection{Calibration guarantee for the discrete-hazard head}

\begin{proposition}[Calibration of the discrete-hazard CDF]
\label{prop:calibration}
Let $\hat\lambda_k=\sigma(g_\phi(H_t,\Delta t_k))$ be the discrete-hazard
output of a finetuned head trained to minimise binary cross-entropy
against per-horizon labels $E_{t+\Delta t_k}$. Define the cumulative
CDF $\hat F_K(H_t)=1-\prod_{k=1}^K(1-\hat\lambda_k)$ and the true CDF
$F_K(H_t)=\PP(E_{t+\Delta t_K}=1\mid H_t)$. Suppose:
\begin{enumerate}[label=(C\arabic*),leftmargin=2em]
\item the head is sufficiently expressive that BCE convergence implies
  pointwise convergence of $\hat\lambda_k(h)\to\lambda_k(h)\coloneqq
  \PP(E_{t+\Delta t_k}=1\mid E_{t+\Delta t_{k-1}}=0,H_t=h)$;
\item the per-horizon hazard rates are bounded:
  $\lambda_k(h)\in[\lambda_{\min},\lambda_{\max}]$ with $0<\lambda_{\min}\le
  \lambda_{\max}<1$.
\end{enumerate}
Then for every $h$,
\begin{equation}\label{eq:calib}
|\hat F_K(h)-F_K(h)|\;\le\;\sum_{k=1}^K\frac{|\hat\lambda_k(h)-\lambda_k(h)|}{1-\lambda_{\max}}.
\end{equation}
\end{proposition}

\begin{proof}
By the telescoping identity
$F_K=1-\prod_{k=1}^K(1-\lambda_k)$ and a first-order Taylor expansion of
$\prod(1-\lambda_k)$ around $(\lambda_1,\dots,\lambda_K)$, the partial
derivative w.r.t.\ $\lambda_j$ has magnitude
$\prod_{k\ne j}(1-\lambda_k)\le 1$. The denominator $1-\lambda_{\max}$
absorbs the worst-case partial.
\end{proof}

\textbf{Implication.} If the per-horizon BCE training error per head is
$O(1/\sqrt n)$ (standard rate), then the cumulative CDF calibration error
is $O(K/\sqrt n)$. For $K=200$ horizons on MBA at $n\approx 10^5$
windows, the calibration bound predicts
$\sim 200/\sqrt{10^5}\approx 0.6$, which is loose --- the bound is
\emph{conservative} (sums independent errors with no cancellation).
Tightening to a $\sqrt{K}/\sqrt n$ rate via a martingale argument is
plausible but is left as future work.

\subsection{Architecture design rules}

We compile the formal prescriptions of Phase 6b plus the new results of
Phase 6c into a single list. Each rule is tagged with its supporting
result.

\begin{enumerate}[leftmargin=1.5em,label=(R\arabic*)]
\item \textbf{Use bidirectional attention pooling for the target encoder.}
A1 requires $H^*$ to be a sufficient statistic for $E$; bidirectional
attention is universal under mild conditions. \emph{Source:}
\S\ref{sec:phase6b} (P2).

\item \textbf{Choose normalisation per dataset family.}
RevIN violates A1 on lifecycle datasets where the drift IS the precursor.
Use train-set global z-score on degradation data (C-MAPSS) and per-instance
RevIN on local-anomaly streams (SMAP/MSL/PSM/SMD/MBA). \emph{Source:}
\S\ref{sec:phase6b} (P4); empirical confirmation v27 (FD001 +0.308 Δt=150 AUROC).

\item \textbf{Sample horizons proportional to $I(H^*;E_{t+\Delta t})/\varepsilon(\Delta t)$.}
The instantaneous information rate. Predicts dense short-Δt sampling on
anomaly data (consistent with v23 patch-tokenisation gains on SMAP) and
near-uniform sampling on degradation data. \emph{Source:} \S\ref{sec:phase6b}
(P3) and Proposition~\ref{prop:per-horizon}.

\item \textbf{Use a discrete-hazard output head, not MonotoneCDF.}
The hazard parameterisation respects the non-monotone per-horizon
posterior $\PP(E_{t+\Delta t}\mid H_t)$ that the FAM labelling
convention generates. \emph{Source:} \S\ref{sec:monotone},
Proposition~\ref{prop:monotone}; empirical confirmation v30 phase 0
(MonotoneCDF h-AUROC $=0.500$).

\item \textbf{Set $d\gtrsim I(H^*;E)/\log_2 d$.}
Rate-distortion bound on representation capacity. For our benchmark,
$d=256$ has substantial headroom over the worst-case requirement.
\emph{Source:} \S\ref{sec:phase6b} (P1).

\item \textbf{Do not transfer predictor weights, only architecture.}
Proposition~\ref{prop:codomain}: pretrained predictor weights have
$\le\sqrt{K/d}$ alignment with finetuning weights in the most favourable
case. \emph{Source:} \S\ref{sec:phase6b}; empirical confirmation v24
phase-12 ablation (RESULTS.md, ``predictor pretraining adds at most a
small boost''). The v30 phase-1 evidence that FAM-predft on FD001 is
within seed noise of FAM-mlp-rand ($0.714\pm 0.028$ vs.\ $0.707\pm 0.018$,
file \texttt{phase1\_decision.json}) further supports this.

\item \textbf{Pretraining provides label-efficiency advantage only when
$\mathfrak R_n(\mathcal H)$ exceeds the representation gap.}
Theorem~\ref{thm:excess}. At 5\% labels on FD001 ($n\approx 5\cdot 10^3$),
$\mathfrak R_n\sim 0.014$ comparable to the gap; pretraining helps
(v20 pred-FT $0.261$ vs.\ scratch $0.035$ F1w). At 10\% labels
($n\approx 10^4$), $\mathfrak R_n\sim 0.010$ smaller than the gap; pred-FT
$\approx$ MLP-rand-init (v30 phase-1: $0.733\pm 0.042$ vs.\
$0.734\pm 0.021$ on FD001).
\end{enumerate}

% =====================================================================
\section{Honest gaps we could not close}

\begin{enumerate}[leftmargin=1.5em]
\item \WEAK\ \textbf{Quantitative $I(H^*;E)$ on each dataset.}
Phase 6b's mapping is qualitative. We do not have a numerical estimate
of $I(H^*;E)$ per dataset; doing so requires either MINE-style estimation
(\citealt{belghazi2018mine}) or a parametric MI estimator on the trained
$H^*$. We list this as the most useful next experiment to make
Proposition 1 quantitatively predictive.

\item \WEAK\ \textbf{The constant $\widetilde C_p$ is data-dependent.}
The fix in \S\ref{sec:phase6a} replaces $C_p$ with $\widetilde C_p$ in
terms of the support of $\eta(H^*)$, which is itself a learned object.
A worst-case version using $\eta(h)\in[0,1]$ trivially gives
$\widetilde C_p=2$ but loses tightness. A proof that
$\widetilde C_p$ is bounded by a function of A2 and A3 alone is still
open.

\item \WEAK\ \textbf{Proposition~\ref{prop:codomain} is for linear heads.}
The full proof in \S\ref{appx:codomain-proof} relies on linearisation.
Extending to nonlinear MLP heads requires careful treatment of the
parameter geometry; we conjecture the bound holds with a Lipschitz
constant prefactor but have not proved it.

\item \WEAK\ \textbf{Calibration bound is loose.}
Proposition~\ref{prop:calibration} gives an $O(K/\sqrt n)$ rate; we
believe the true rate is $O(\sqrt{K/n})$ via a martingale argument but
have not formalised it.

\item \WEAK\ \textbf{Per-dataset $L$ is unmeasured.}
The Lipschitz constant $L$ in A3 is treated as a fixed but unknown
parameter. Estimating it empirically would let us produce
quantitative versions of the predictions in Phase 6b.

\item \WEAK\ \textbf{BATADAL ambiguity.}
On BATADAL, FAM-probe ($0.521\pm 0.038$) is at chance but FAM-predft
($0.607\pm 0.033$) is well above. Proposition 1 predicts ``moderate
$I(H_t;E)$ implies moderate downstream'' but cannot disambiguate (a)
$I(H_t;E)$ is in fact high but the linear probe cannot extract it
versus (b) $I(H_t;E)$ is low and the head's nonlinearity is
hallucinating. Resolving this would require a finer
information-vs-extractability decomposition (cf.\ ``usable information''
of \citealp{xu2020theory}).
\end{enumerate}

% =====================================================================
\section{Appendix: full proof of Proposition~\ref{prop:codomain}}
\label{appx:codomain-proof}

\begin{proof}[Proof of Proposition~\ref{prop:codomain} (linear case)]
Let $\predictor(h,\Delta t;\phi)=W h$ for a parameter matrix
$W\in\R^{d\times d}$ in the pretraining case and
$W\in\R^{K\times d}$ in the finetuning case (we drop the $\Delta t$
dependence by absorbing it into the input). The pretraining objective is
$\E\|W H_t-H^*\|_2^2$ minimised by
$W^{\star}_{\mathrm{pre}}=\Sigma_{H^*H_t}\Sigma_{H_tH_t}^{-1}$, the
linear regression of $H^*$ onto $H_t$.

Let $\Pi_E\in\R^{d\times d}$ be the orthogonal projector onto the
event-relevant subspace, defined as the row-span of the linear regression
of $E$ on $H^*$ (a one-dimensional subspace since $E$ is binary, hence
$\mathrm{rank}(\Pi_E)=1$). Then
\[
\|\Pi_E W^{\star}_{\mathrm{pre}}\|_F^2
=\Pi_E^{\top}\Sigma_{H^*H_t}\Sigma_{H_tH_t}^{-1}\Sigma_{H_tH^*}\Pi_E.
\]
The denominator
$\|W^{\star}_{\mathrm{pre}}\|_F^2=\mathrm{tr}(\Sigma_{H^*H_t}
\Sigma_{H_tH_t}^{-1}\Sigma_{H_tH^*})$ is the total
predictive trace.

By Gaussian-MI duality (Cover \& Thomas, Ch.~8), for jointly Gaussian
$(H^*,H_t)$:
$I(H_t;H^*)=\tfrac12\log\det(I+\Sigma_{H^*H_t}\Sigma_{H_tH_t}^{-1}\Sigma_{H_tH^*}\Sigma_{H^*H^*}^{-1})$;
likewise $I(H^*;E)=\tfrac12\log(1+\mathrm{snr}(E\mid H^*))$ in
linearised form. Taking ratios of the dominant eigenvalues gives
\eqref{eq:codomain-bound} after some algebra. The non-Gaussian case
follows by the standard linearisation argument with an extra factor for
the curvature of the cross-entropy loss; we have not made this fully
rigorous.
\end{proof}

% =====================================================================
\section{Summary}

\begin{itemize}[leftmargin=1.5em]
\item \CONF\ All seven proof steps of Proposition 1 are correct under
the stated assumptions, \emph{except} that the Jensen-gap step requires
a calibrated-posterior assumption A1$'$ that is not currently numbered in
\texttt{theory\_main.tex}. We documented the fix without editing the
in-paper proof.
\item \CONF\ The qualitative prediction of Proposition 1 + Corollary 2
matches 12/13 datasets in RESULTS.md (the remaining one is
inconclusive, $n=1$ seed).
\item \CONF\ Proposition~\ref{prop:codomain} formalises the
codomain-mismatch argument and predicts the v24/v30 finding that
predictor pretrained weights do not transfer.
\item \CONF\ Theorem~\ref{thm:excess} explains both the v20 5\%-labels
result (pretraining helps) and the v30 phase-1 10\%-labels result
(pretraining $\approx$ random init), via a representation-gap vs.\
statistical-error trade-off.
\item \CONF\ Proposition~\ref{prop:per-horizon} predicts the observed
declining per-horizon AUROC tail.
\item \CONF\ Proposition~\ref{prop:monotone} explains why the MonotoneCDF
head failed in v30 phase 0 ($0.500$ vs.\ $0.813$ for the discrete-hazard
head): FAM's per-horizon labelling generates a non-monotone posterior
that a monotone head cannot fit.
\item \OPEN\ Quantitative estimates of $I(H^*;E)$ and $L$ per dataset.
\item \OPEN\ Tightening Proposition~\ref{prop:calibration} from
$O(K/\sqrt n)$ to $O(\sqrt{K/n})$.
\item \OPEN\ Disambiguating BATADAL via a usable-information analysis.
\end{itemize}

\bibliographystyle{plainnat}
% Inline bibliography to keep this file standalone.
\begin{thebibliography}{99}
\bibitem[Cover \& Thomas(2006)]{cover2006elements}
T.~M. Cover and J.~A. Thomas. \emph{Elements of Information Theory}, 2nd ed. Wiley, 2006.

\bibitem[Tishby et al.(2000)]{tishby2000information}
N.~Tishby, F.~C. Pereira, and W.~Bialek. The information bottleneck method. \emph{Proc. 37th Allerton Conf.}, 2000.

\bibitem[Tosh et al.(2021)]{tosh2021contrastive}
C.~Tosh, A.~Krishnamurthy, and D.~Hsu. Contrastive estimation reveals topic posterior information to linear models. \emph{JMLR}, 22(281), 2021.

\bibitem[Bialek et al.(2001)]{Bialek2001}
W.~Bialek, I.~Nemenman, and N.~Tishby. Predictability, complexity, and learning. \emph{Neural Computation}, 13(11), 2001.

\bibitem[Arora et al.(2019)]{Arora2019contrastive}
S.~Arora, H.~Khandeparkar, M.~Khodak, O.~Plevrakis, and N.~Saunshi. A theoretical analysis of contrastive unsupervised representation learning. \emph{ICML}, 2019.

\bibitem[Liao \& Berg(2019)]{liao2019sharpening}
J.~G. Liao and A.~Berg. Sharpening Jensen's inequality. \emph{The American Statistician}, 73(3), 2019.

\bibitem[Lee et al.(2018)]{lee2018deephit}
C.~Lee, W.~Zame, J.~Yoon, and M.~van~der~Schaar. DeepHit: A deep learning approach to survival analysis with competing risks. \emph{AAAI}, 2018.

\bibitem[Belghazi et al.(2018)]{belghazi2018mine}
M.~I. Belghazi et al. Mutual information neural estimation. \emph{ICML}, 2018.

\bibitem[Yun et al.(2020)]{yun2020universal}
C.~Yun, S.~Bhojanapalli, A.~S.~Rawat, S.~J.~Reddi, and S.~Kumar. Are transformers universal approximators of sequence-to-sequence functions? \emph{ICLR}, 2020.

\bibitem[Mohri et al.(2018)]{MohriRostamizadeh2018}
M.~Mohri, A.~Rostamizadeh, and A.~Talwalkar. \emph{Foundations of Machine Learning}, 2nd ed. MIT Press, 2018.

\bibitem[McShane(1934)]{McShane1934}
E.~J. McShane. Extension of range of functions. \emph{Bull. AMS}, 40(12), 1934.

\bibitem[Durrett(2019)]{Durrett2019}
R.~Durrett. \emph{Probability: Theory and Examples}, 5th ed. Cambridge, 2019.

\bibitem[Xu et al.(2020)]{xu2020theory}
Y.~Xu, S.~Zhao, J.~Song, R.~Stewart, and S.~Ermon. A theory of usable information under computational constraints. \emph{ICLR}, 2020.

\end{thebibliography}

\end{document}
